\section{The day-ahead market and our 09:00 cutoff}

\subsection{The daily timeline (Poland, CET)}

Same rhythm every day, weekends included:

\begin{itemize}
  \item \textbf{Morning, D-1.} Desks update forecasts: load, wind, solar, prices.
        Traders build bids for every hour of tomorrow (day D).
  \item \textbf{12:00, D-1.} Day-ahead auction closes (TGE locally; PL is coupled
        with the EU market via SDAC --- one algorithm, EUPHEMIA, clears Europe).
  \item \textbf{Early afternoon, D-1.} Clearing prices and volumes publish.
        One price per hour for each bidding zone.
  \item \textbf{Day D.} Deviations from the traded schedule settle on the
        balancing market, usually at worse prices. Forecast errors cost money here.
\end{itemize}

\subsection{Why we cut off at 09:00 on D-1}

Our rule: the forecast for day D uses only data that existed at 09:00 on D-1.

\begin{itemize}
  \item The auction closes at 12:00. A forecast delivered later is useless.
  \item Three hours of margin covers data delays and rework --- desk reality.
  \item Every backtest applies the same cutoff. Otherwise backtest accuracy
        is fantasy (see leakage note~03).
\end{itemize}

\subsection{What is known at the cutoff, for day D}

\begin{center}
\begin{tabular}{lp{7cm}}
\toprule
Known & Load history to $\sim$08:00 D-1; weather \emph{forecasts} for D;
        calendar (holidays, weekday) --- known years ahead \\
\midrule
Unknown & Actual load of D-1 afternoon and D; actual weather of D;
          the auction price of D \\
\bottomrule
\end{tabular}
\end{center}

Consequence: the freshest usable ``same hour'' load lag is 48~hours
(day D-2), not 24. Lag~24 would need data that arrives only after the cutoff.

\subsection{Interview line}

``I time-stamp every input against a 09:00 D-1 cutoff because that is when a
day-ahead desk actually commits. My backtest cannot use anything the desk
would not have had.''
